\documentclass[11pt,a4paper]{article}

% --- PAQUETS DE BASE ---
\usepackage[utf8]{inputenc}
\usepackage[T1]{fontenc}
\usepackage[french]{babel}
\usepackage{amsmath, amsfonts, amssymb}
\usepackage{geometry}
\geometry{top=2cm, bottom=2.5cm, left=2cm, right=2cm}

% --- DESIGN ---
\usepackage{xcolor}
\usepackage{colortbl}
\definecolor{BrandPrimary}{HTML}{0D9488}
\definecolor{BrandDark}{HTML}{111827}

\usepackage{tcolorbox}
\tcbuselibrary{skins, breakable}

\usepackage{titlesec}
\titleformat{\section}{\color{BrandPrimary}\normalfont\Large\bfseries}{\thesection}{1em}{}[{\titlerule[0.8pt]}]

% --- POLICE ---
\usepackage{helvet}
\renewcommand{\familydefault}{\sfdefault}

\begin{document}

\begin{tcolorbox}[enhanced, colback=BrandDark, colframe=BrandPrimary, arc=10pt, boxrule=2pt, halign=center]
    \vspace{0.3cm}
    {\Huge \bfseries \color{white} ZENITH LEARN} \\
    \vspace{0.2cm}
    {\Large \color{BrandPrimary} IA Éthique : Modélisation TWIST 05} \\
    \vspace{0.3cm}
\end{tcolorbox}

\section{Analyse sémantique : La Loi des 20/80}
Le \textbf{TWIST 05} établit une asymétrie de population : \textit{``20\% des utilisateurs sont déjà entrepreneurs actifs, 80\% explorent encore''}. Cela interdit un modèle unique de recommandation ``One size fits all''.

Cette répartition, familière sous le nom de \textbf{Principe de Pareto}, reflète une réalité psychologique profonde : la plupart des individus sur une plateforme d'apprentissage entrepreneurial sont en phase de \textit{découverte}. Ils cherchent à comprendre si l'entrepreneuriat est fait pour eux, à identifier leur domaine, à construire leur confiance. Une minorité, déjà engagée dans un projet concret, a des besoins radicalement différents : elle cherche des solutions précises à des problèmes immédiats. Recommander la même chose aux deux groupes serait aussi absurde que de proposer le même régime alimentaire à un marathonien et à un sédentaire.

\section{Modélisation de la Segmentation Dynamique}
L'IA doit s'adapter en modifiant ses vecteurs de pondération en fonction du profil identifié.

\subsection{Formule du Score de Recommandation Segmenté}
\begin{equation}
Score(u, c) = \vec{W}_{Seg(u)} \cdot \vec{F}_c
\end{equation}

\subsection{Genèse de la formule : pourquoi cette construction ?}

L'écriture vectorielle de cette formule est fondamentale. Elle exprime que le score d'un cours $c$ pour un utilisateur $u$ est un \textbf{produit scalaire} entre :
\begin{itemize}
    \item $\vec{W}_{Seg(u)}$ : le vecteur de poids \textit{propre au segment} de l'utilisateur. Ce vecteur change selon qu'on est Explorateur ou Acteur.
    \item $\vec{F}_c$ : le vecteur des features du cours $c$ (son score d'engagement, sa complétion ajustée, son signal d'action, etc.).
\end{itemize}

Le produit scalaire est la somme des produits terme à terme : chaque feature du cours est multipliée par le poids que le \textit{segment de l'utilisateur} lui accorde. La puissance de cette écriture est de permettre de changer radicalement l'évaluation d'un même cours selon le profil de l'apprenant, simplement en changeant le vecteur $\vec{W}$.

\subsection{Détail et Justification Mathématique du Profilage}
\begin{itemize}
    \item \textbf{$Seg(u) \in \{Active, Explorer\}$ (Variable Latente) :}
    Le segment est déterminé par un classifieur binaire basé sur le ratio Action/Engagement (T04). Concrètement, si $SPI(u)$ est dominé par ses composantes terrain ($\bar{A_u}$, $B_u$), l'utilisateur est classé \textit{Acteur}. Si $SPI(u)$ est principalement alimenté par des signaux académiques, il est classé \textit{Explorateur}. Ce classifieur est dit ``latent'' car le segment n'est pas renseigné par l'utilisateur lui-même~: il est \textit{inféré} par le modèle à partir du comportement.

    \item \textbf{$\vec{W}_{Explorer} = [0.5, 0.5, 0, 0]$ (Vecteur Explorateur) :}
    Pour les 80\% d'explorateurs, on privilégie la \textbf{Découverte} en donnant 50\% du poids à l'engagement ($S_E$) et 50\% à la complétion ajustée ($S_C^{adj}$). Les termes liés à l'action terrain (colonnes 3 et 4) sont nuls, car ces apprenants n'ont pas encore de projet : leur recommander un cours très orienté ``lancement d'entreprise'' serait prématuré et potentiellement décourageant. L'IA les guide vers les contenus de découverte qui ont le mieux fonctionné chez des explorateurs similaires.

    \item \textbf{$\vec{W}_{Active} = [0.1, 0.1, 0.4, 0.4]$ (Vecteur Acteur) :}
    Pour les 20\% d'acteurs, on privilégie l'\textbf{Efficacité}. L'action terrain et le lancement d'entreprise dominent avec un poids cumulé de 80\%. L'IA cherche ici le cours qui débloquera la prochaine étape concrète de leur business. L'engagement (0.1) et la complétion (0.1) ne servent plus que de filtre qualité minimal : on s'assure que le cours recommandé est bien reçu, mais l'IA ne cherche plus à maximiser le plaisir de la découverte, elle cherche à maximiser l'impact opérationnel.

    \item \textbf{Dynamisme de la segmentation :} Le segment $Seg(u)$ n'est pas figé. Un explorateur qui accumule suffisamment de signaux d'action peut basculer dans le segment Acteur. Ce dynamisme permet à l'IA d'accompagner la progression naturelle de l'entrepreneur, en adaptant ses recommandations au moment précis du parcours.
\end{itemize}

\section{Nouvelle Structure du Dataset (Spécification T05)}
\begin{table}[h!]
\centering
\begin{tabular}{|l|l|p{7cm}|}
\hline
\rowcolor{BrandDark} \color{white} Colonne & \rowcolor{BrandDark} \color{white} Type & \rowcolor{BrandDark} \color{white} Rôle T05 \\ \hline
\rowcolor{BrandPrimary!20} \texttt{user\_segment} & Enum & $[\text{Explorateur, Actif}]$. Segment inféré. \\ \hline
\texttt{is\_explorer} & Binary & Utilisé pour injecter le vecteur de pondération correspondant dans le calcul du score. \\ \hline
\end{tabular}
\end{table}

\section{Impact sur les acquis précédents}
\begin{itemize}
    \item \textbf{Complexification du TWIST 04} : Le score de potentiel entrepreneurial ($SPI$) n'est plus une sortie statique, mais une entrée de segmentation. Il sert à classer l'utilisateur, ce qui détermine ensuite quels poids seront appliqués dans les recommandations futures. Le $SPI$ devient ainsi une variable à double rôle : score de potentiel \textit{et} critère de segmentation.
    \item \textbf{Affinement du TWIST 02/03} : Les règles de sobriété et de redressement de durée s'appliquent de manière différenciée. Un entrepreneur actif tolère moins de ``bruit'' pédagogique qu'un explorateur : chaque heure passée sur une formation a un coût d'opportunité bien réel pour lui. Les corrections T02 et T03 sont donc \textit{plus importantes} pour le segment Acteur.
\end{itemize}

\section{Conclusion}
Le TWIST 05 transforme Zenith Learn d'un catalogue uniforme en une plateforme capable de mener de front deux missions : l'\textbf{éveil entrepreneurial} (80\%) et l'\textbf{accélération d'impact} (20\%). Il introduit le concept de \textbf{personnalisation contextuelle} : ce n'est plus seulement ``quel cours correspond à cet utilisateur ?'', mais ``quel cours correspond à cet utilisateur \textit{là où il en est de son parcours entrepreneurial} ?''.

\end{document}